\documentclass[12pt]{article}
\usepackage[a4paper,margin=1in]{geometry}
\usepackage{amsmath,amssymb}
\usepackage{enumitem}
\usepackage[a4paper,margin=1in]{geometry}
\usepackage{amsmath,amssymb}
\usepackage{mathtools}
\usepackage{hyperref}
\usepackage{titlesec}
\title{Main results - MASA03}
\author{}
\date{}
\begin{document}
\maketitle
\section*{Probability spaces}
Definition. A \emph{probability space} is a triple \((\Omega,\mathcal F,P)\) where:
\begin{itemize}[leftmargin=1.2em]
  \item \(\Omega\): \emph{sample space} (set of outcomes).
  \item \(\mathcal F\subseteq 2^\Omega\): \emph{$\sigma$-algebra} of events: 
    \(\Omega\in\mathcal F\); if \(A\in\mathcal F\) then \(A^c\in\mathcal F\); 
    if \(A_1,A_2,\dots\in\mathcal F\) then \(\bigcup_{n\ge1}A_n\in\mathcal F\).
  \item \(P:\mathcal F\to[0,1]\): \emph{probability measure} satisfying the axioms:
  \begin{enumerate}[label=(\roman*), leftmargin=1.5em]
    \item Non-negativity: \(P(A)\ge 0\) for all \(A\in\mathcal F\).
    \item Normalization: \(P(\Omega)=1\).
    \item $\sigma$-additivity: if \(A_i\in\mathcal F\) are pairwise disjoint,
      \(\displaystyle P\!\Big(\bigcup_{i=1}^\infty A_i\Big)=\sum_{i=1}^\infty P(A_i)\).
  \end{enumerate}
\end{itemize}

Immediate consequences (useful):
\(P(\varnothing)=0\), \quad \(P(A^c)=1-P(A)\), \quad 
\(A\subseteq B\Rightarrow P(A)\le P(B)\), \quad
if \(A\cap B=\varnothing\) then \(P(A\cup B)=P(A)+P(B)\).


\subsection{Random Variable}
Definition. A function $X: \Omega \to \mathbb{R}$ is a \emph{random variable} on the probability space $(\Omega, \mathcal F, P)$ if it satisfies the condition of $\mathcal F$-measurability:
$$
\text{For every Borel set } B \in \mathcal{B}(\mathbb{R}), \quad X^{-1}(B) = \{\omega \in \Omega : X(\omega) \in B\} \in \mathcal{F}
$$
\subsubsection*{Required Conditions}
The crucial condition is that the pre-image of any set $B$ for which a probability is defined in $\mathbb{R}$ must be an event in the $\sigma$-algebra $\mathcal{F}$. The probability measure $P$ on $\mathcal{F}$ itself is not involved in defining $X$ as a random variable, only the structure of $\mathcal{F}$.

\section*{Law of Total Probability (LTP)}
If \(A_1,\dots,A_k\) form a partition of \(\Omega\) (pairwise disjoint, \(\bigcup_i A_i=\Omega\)), then for any event \(B\),
\[
P(B)=\sum_{i=1}^k P(B\mid A_i)\,P(A_i).
\]
(Countable version holds similarly.)

\section*{Bayes' Rule}
For events \(A_j\) and \(B\) with \(P(B)>0\),
\[
P(A_j\mid B)=\frac{P(B\mid A_j)\,P(A_j)}{\sum_{i} P(B\mid A_i)\,P(A_i)}.
\]
The denominator is \(P(B)\) from the LTP over the partition \(\{A_i\}\).

\section{Transforms of random variables}

\subsection{Sums (convolution)}
\paragraph{Continuous, independent:}
If $Z=X+Y$ and $X\perp Y$ with densities $f_X,f_Y$,
\[
f_Z(z)=(f_X * f_Y)(z)=\int_{-\infty}^{\infty} f_X(x)\,f_Y(z-x)\,dx.
\]
\textit{Support rule:} integrate $x$ over the intersection of supports
$\{x:\ f_X(x)>0\}$ and $\{x:\ f_Y(z-x)>0\}$; this often yields piecewise limits.

\paragraph{Discrete, independent:}
If $p_X,p_Y$ are pmfs and $Z=X+Y$,
\[
p_Z(z)=\sum_{k\in\mathbb Z} p_X(k)\,p_Y(z-k).
\]

\paragraph{General (possibly dependent):}
If the joint density $f_{X,Y}$ is known,
\[
f_Z(z)=\int_{-\infty}^{\infty} f_{X,Y}\big(x,\,z-x\big)\,dx.
\]

\subsection{Monotone scalar transforms}
If $Z=g(X)$ and $g$ is strictly monotone with inverse $g^{-1}$,
\[
f_Z(z)= f_X\!\big(g^{-1}(z)\big)\;\left|\frac{d}{dz}g^{-1}(z)\right|.
\]
(\textit{CDF method:} $F_Z(z)=P(g(X)\le z)$ also works when $g$ is not monotone.)

\subsection{Two-to-two transforms (Jacobian)}
If $(U,V)=T(X,Y)$ is a bijection with inverse $(X,Y)=T^{-1}(U,V)$,
\[
f_{U,V}(u,v)= f_{X,Y}\!\big(T^{-1}(u,v)\big)\;
\left|\det \frac{\partial(x,y)}{\partial(u,v)}\right|.
\]
Marginals follow by integrating out the other variable.

\subsection{Products and quotients}
\paragraph{Product $Z=XY$ (general case):}
\[
f_Z(z)=\int_{-\infty}^{\infty} f_{X,Y}\!\Big(x,\frac{z}{x}\Big)\,\frac{1}{|x|}\,dx
\quad (\text{integrate over } x \text{ where } (x,z/x) \text{ is in support}).
\]
If $X,Y$ are independent: replace $f_{X,Y}$ with $f_X f_Y$ inside the integral.

\paragraph{Quotient $Z=X/Y$:}
\[
f_Z(z)=\int_{-\infty}^{\infty} f_{X,Y}\!\big(zy,y\big)\,|y|\,dy.
\]

\subsection{Min and Max (use CDFs)}
Let $M=\max(X,Y)$ and $m=\min(X,Y)$.
If $X,Y$ are independent with cdfs $F_X,F_Y$ and densities $f_X,f_Y$,
\[
F_M(t)=P(M\le t)=F_X(t)F_Y(t),\qquad
f_M(t)= f_X(t)F_Y(t)+F_X(t)f_Y(t).
\]
\[
F_m(t)=1-P(X>t\ \text{or}\ Y>t)=1-(1-F_X(t))(1-F_Y(t)),\quad
f_m(t)= f_X(t)(1-F_Y(t))+f_Y(t)(1-F_X(t)).
\]

\subsection{Quick recipe}
\begin{enumerate}[leftmargin=1.2em]
  \item \textbf{Identify the map:} $Z=X+Y$, $Z=g(X)$, or $(U,V)=T(X,Y)$.
  \item \textbf{Pick the tool:} convolution (sum), monotone transform (one-to-one in 1D),
        Jacobian (two-to-two), or the CDF method.
  \item \textbf{Set the limits by geometry:} intersect supports after the change of variables.
  \item \textbf{Normalize/check:} verify $\int f_Z=1$ (or $\sum p_Z=1$).
\end{enumerate}

\section{Joint distributions and tools}

\subsection{Joint density and normalization}
A pair $(X,Y)$ has joint density $f_{X,Y}(x,y)\ge 0$ with
\[
\iint_{\mathbb R^2} f_{X,Y}(x,y)\,dx\,dy=1.
\]

\subsection{Marginals}
\[
f_X(x)=\int_{-\infty}^{\infty} f_{X,Y}(x,y)\,dy,
\qquad
f_Y(y)=\int_{-\infty}^{\infty} f_{X,Y}(x,y)\,dx.
\]

\subsection{Independence}
$X$ and $Y$ are independent iff
\[
f_{X,Y}(x,y)=f_X(x)\,f_Y(y)\quad\text{(almost everywhere on the support).}
\]

\subsection{Conditional density}
For $f_X(x)>0$,
\[
f_{Y\mid X}(y\mid x)=\frac{f_{X,Y}(x,y)}{f_X(x)}.
\]

\subsection{Cumulative distribution function (cdf)}
For any random variable $Z$,
\[
F_Z(z)=\Pr(Z\le z), \qquad \Pr(a<Z\le b)=F_Z(b)-F_Z(a).
\]
If $Z$ has a density $f_Z$, then
\[
f_Z(z)=\frac{d}{dz}F_Z(z).
\]

\subsection{Product variable via the cdf method}
For $Z=XY$ with $(X,Y)$ supported on $[a,b]\times[c,d]\subseteq[0,\infty)^2$,
\[
F_Z(z)=\Pr(XY\le z)=
\iint_{\{(x,y)\in[a,b]\times[c,d]:\,xy\le z\}} f_{X,Y}(x,y)\,dx\,dy,
\]
and $f_Z(z)=\dfrac{d}{dz}F_Z(z)$ on the interior of the support of $Z$.

\section{Expectation and related results}

\subsection{Definition}
For a random variable $X$ with pdf $f_X$ or pmf $p_X$,
\[
E[X]=
\begin{cases}
\displaystyle \sum_x x\,p_X(x), & \text{(discrete)},\\[2mm]
\displaystyle \int_{-\infty}^{\infty} x\,f_X(x)\,dx, & \text{(continuous)}.
\end{cases}
\]
More generally (Law of the Unconscious Statistician, LOTUS),
\[
E[g(X)] =
\begin{cases}
\displaystyle \sum_x g(x)\,p_X(x),\\[2mm]
\displaystyle \int g(x)\,f_X(x)\,dx.
\end{cases}
\]

\subsection{Conditional expectation and total laws}
For densities,
\[
E[X\mid Y=y]=\int x\,f_{X\mid Y}(x\mid y)\,dx,
\qquad
f_{X\mid Y}(x\mid y)=\frac{f_{X,Y}(x,y)}{f_Y(y)}.
\]

\paragraph{Law of total expectation (tower property):}
\[
E[X]=E\big(E[X\mid Y]\big).
\]

\paragraph{Law of total probability (continuous form):}
For any event $A$,
\[
P(A)=\int P(A\mid Y=y)\,f_Y(y)\,dy
\quad\text{or, for discrete $Y$,}\quad
P(A)=\sum_y P(A\mid Y=y)\,P(Y=y).
\]

This expresses the overall probability $P(A)$ as the average (expectation) of the conditional probabilities $P(A\mid Y=y)$ weighted by how likely each $y$ is.  
Equivalently,
\[
P(A)=E\big[P(A\mid Y)\big].
\]
Hence, probabilities are special cases of expectations:
\[
P(X<Y)=E\big[P(X<Y\mid X)\big].
\]

\subsection{Variance and Covariance}
\begin{align*}
\operatorname{Var}(X) &= E[X^2]-\big(E[X]\big)^2 \\
\operatorname{Cov}(X,Y) &= E[XY]-E[X]E[Y]
\end{align*}

\subsubsection{Law of Total Variance (Conditional Version)}
\[
\operatorname{Var}(X)=E[\operatorname{Var}(X\mid Y)] + \operatorname{Var}(E[X\mid Y])
\]

\subsubsection{Relationship Between $\operatorname{Cov}$ and $\operatorname{Var}$}
\[
\operatorname{Var}(X) = \operatorname{Cov}(X, X)
\]

\subsubsection{Correlation Coefficient ($\rho$)}
The correlation coefficient is the standardized covariance.
\[
\rho_{X,Y} = \operatorname{Corr}(X,Y) = \frac{\operatorname{Cov}(X,Y)}{\sqrt{\operatorname{Var}(X)}\sqrt{\operatorname{Var}(Y)}} = \frac{\operatorname{Cov}(X,Y)}{\sigma_X \sigma_Y}
\]
\paragraph{Identity:}
The identity for correlation is its range:
\[
-1 \leq \rho_{X,Y} \leq 1
\]
\paragraph{Covariance in terms of Correlation:}
\[
\operatorname{Cov}(X,Y) = \rho_{X,Y} \sigma_X \sigma_Y
\]

\section{Common Probability Distributions}

\subsection{Discrete distributions}

\subsubsection*{Binomial distribution}
\[
X \sim \operatorname{Bin}(n,p)
\]
\[
P(X = k) = \binom{n}{k} p^k (1-p)^{n-k}, \qquad k = 0,1,\dots,n.
\]
\[
E[X] = np, \qquad \operatorname{Var}(X) = np(1-p).
\]

\subsubsection*{Bernoulli distribution}
\[
X \sim \operatorname{Bern}(p)
\]
\[
P(X = 1) = p, \quad P(X = 0) = 1 - p.
\]
\[
E[X] = p, \qquad \operatorname{Var}(X) = p(1-p).
\]

\subsubsection*{Geometric distribution}
Two common parameterizations with success probability $p\in(0,1)$:

\paragraph{(Trials until first success) } 
$X \sim \operatorname{Geom}_1(p)$ with support $\{1,2,\ldots\}$.
\[
P(X=k)=(1-p)^{k-1}p,\qquad k=1,2,\ldots
\]
\[
E[X]=\frac{1}{p},\qquad \operatorname{Var}(X)=\frac{1-p}{p^{2}}.
\]
Memoryless: $P(X>m+n\mid X>m)=(1-p)^n$.

\paragraph{(Failures before first success) } 
$Y \sim \operatorname{Geom}_0(p)$ with support $\{0,1,2,\ldots\}$.
\[
P(Y=k)=(1-p)^{k}p,\qquad k=0,1,2,\ldots
\]
\[
E[Y]=\frac{1-p}{p},\qquad \operatorname{Var}(Y)=\frac{1-p}{p^{2}}.
\]
Relation: $X=Y+1$.

\paragraph{CDFs:}
\[
F_X(k)=1-(1-p)^{k},\quad k\in\{1,2,\ldots\};\qquad
F_Y(k)=1-(1-p)^{k+1},\quad k\in\{0,1,2,\ldots\}.
\]

\paragraph{Notes:} $\operatorname{Geom}_1(p)$ is the discrete analog of $\operatorname{Exp}(\lambda)$ and is a special case of the negative binomial (number of trials until the $r$-th success with $r=1$).

\subsubsection*{Poisson distribution}
\[
X \sim \operatorname{Poisson}(\lambda)
\]
\[
P(X = k) = \frac{e^{-\lambda} \lambda^k}{k!}, \qquad k = 0,1,2,\dots
\]
\[
E[X] = \lambda, \qquad \operatorname{Var}(X) = \lambda.
\]

\subsection{Continuous distributions}

\subsubsection*{Uniform distribution}
\[
X \sim \operatorname{Unif}(a,b)
\]
\[
f_X(x) = 
\begin{cases}
\dfrac{1}{b - a}, & a \le x \le b,\\
0, & \text{otherwise.}
\end{cases}
\]
\[
E[X] = \dfrac{a + b}{2}, \qquad \operatorname{Var}(X) = \dfrac{(b - a)^2}{12}.
\]

\subsubsection*{Exponential distribution}
\[
X \sim \operatorname{Exp}(\lambda)
\]
\[
f_X(x) = 
\begin{cases}
\lambda e^{-\lambda x}, & x \ge 0,\\
0, & \text{otherwise.}
\end{cases}
\]
\[
E[X] = \dfrac{1}{\lambda}, \qquad \operatorname{Var}(X) = \dfrac{1}{\lambda^2}.
\]
\textbf{Memoryless property:} 
\(P(X > s + t \mid X > s) = P(X > t).\)

\subsubsection*{Normal (Gaussian) distribution}
\[
X \sim N(\mu, \sigma^2)
\]
\[
f_X(x) = \frac{1}{\sqrt{2\pi\sigma^2}}
\exp\!\left(-\frac{(x - \mu)^2}{2\sigma^2}\right), \qquad x \in \mathbb{R}.
\]
\[
E[X] = \mu, \qquad \operatorname{Var}(X) = \sigma^2.
\]

\subsubsection*{Standard Normal distribution}
\[
Z \sim N(0,1)
\]
\[
f_Z(z) = \frac{1}{\sqrt{2\pi}} e^{-z^2/2},
\qquad
F_Z(z) = \Phi(z) = \int_{-\infty}^{z} f_Z(t)\,dt.
\]

\subsubsection*{Gamma distribution}
\[
X \sim \operatorname{Gamma}(\alpha, \lambda)
\]
\[
f_X(x) =
\begin{cases}
\dfrac{\lambda^{\alpha}}{\Gamma(\alpha)}\, x^{\alpha-1} e^{-\lambda x}, & x > 0,\\
0, & \text{otherwise.}
\end{cases}
\]
\[
E[X] = \dfrac{\alpha}{\lambda}, \qquad \operatorname{Var}(X) = \dfrac{\alpha}{\lambda^2}.
\]
Special case: \(\operatorname{Exp}(\lambda) = \operatorname{Gamma}(1,\lambda).\)

\subsection{Summary Table}

\begin{center}
\begin{tabular}{lccc}
\hline
\textbf{Distribution} & \textbf{Support} & \textbf{Mean} & \textbf{Variance} \\ \hline
Bernoulli($p$) & $\{0,1\}$ & $p$ & $p(1-p)$ \\
Binomial($n,p$) & $\{0,\dots,n\}$ & $np$ & $np(1-p)$ \\
Poisson($\lambda$) & $\mathbb{N}_0$ & $\lambda$ & $\lambda$ \\
Uniform($a,b$) & $[a,b]$ & $(a+b)/2$ & $(b-a)^2/12$ \\
Exponential($\lambda$) & $[0,\infty)$ & $1/\lambda$ & $1/\lambda^2$ \\
Normal($\mu,\sigma^2$) & $\mathbb{R}$ & $\mu$ & $\sigma^2$ \\
Gamma($\alpha,\lambda$) & $[0,\infty)$ & $\alpha/\lambda$ & $\alpha/\lambda^2$ \\
\hline
\end{tabular}
\end{center}

\section{Limit Theorems: LLN and CLT}

\subsection{Law of Large Numbers (LLN)}

Let $X_1, X_2, \dots$ be i.i.d.\ random variables with $E[X_i]=\mu$ and $\operatorname{Var}(X_i)=\sigma^2<\infty$.
Define the sample mean
\[
\bar{X}_n = \frac{1}{n}\sum_{i=1}^{n} X_i.
\]

\subsubsection*{Weak Law of Large Numbers (WLLN)}
\[
\bar{X}_n \xrightarrow{P} \mu \quad\text{as } n\to\infty.
\]
That is, for every $\varepsilon>0$,
\[
P(|\bar{X}_n - \mu| > \varepsilon) \longrightarrow 0.
\]
\textbf{Interpretation:} The sample mean converges in probability to the true mean $\mu$.

\subsubsection*{Strong Law of Large Numbers (SLLN)}
\[
\bar{X}_n \xrightarrow{\text{a.s.}} \mu \quad\text{as } n\to\infty.
\]
That is,
\[
P\!\Big(\lim_{n\to\infty}\bar{X}_n = \mu\Big)=1.
\]
\textbf{Interpretation:} The sample mean converges almost surely (with probability 1) to the true mean $\mu$.

\paragraph{Difference between weak and strong forms:}
Both state that $\bar{X}_n$ approaches $\mu$ as $n$ grows, but:
\begin{itemize}[leftmargin=1.2em]
  \item \textbf{Weak LLN:} convergence in probability (approximate closeness for large $n$).
  \item \textbf{Strong LLN:} almost sure convergence (the entire sequence converges for almost every outcome).
\end{itemize}

\subsection{Central Limit Theorem (CLT)}

\textbf{Theorem (Lindeberg-Lévy CLT).}
Let $X_1, X_2, \dots, X_n$ be i.i.d.\ random variables with
\[
E[X_i]=\mu, \qquad \operatorname{Var}(X_i)=\sigma^2<\infty.
\]
Then as $n\to\infty$,
\[
\frac{\sqrt{n}(\bar{X}_n - \mu)}{\sigma}
\xRightarrow{d} N(0,1),
\]
where $\xRightarrow{d}$ denotes convergence in distribution.

\textbf{Interpretation:}  
For large $n$, the sample mean $\bar{X}_n$ is approximately normal:
\[
\bar{X}_n \approx N\!\left(\mu, \frac{\sigma^2}{n}\right).
\]
Hence, the CLT justifies using the normal distribution to approximate sums or averages of many independent random variables.

\paragraph{Practical consequence:}
For large $n$,
\[
P\!\left(\frac{\bar{X}_n - \mu}{\sigma/\sqrt{n}} \le z\right) \approx \Phi(z),
\]
where $\Phi(z)$ is the standard normal cdf.  
This result underpins confidence intervals and hypothesis testing.




\end{document}